%-------------------------------------------------------------------------------
%	SIDEBAR -- DEFINITIONS
%
%	Everything `sidebar` is, apart from its page geometry and the order of the
%	two columns.
%
%	This was once shared with `sidebarleft`, on the assumption that the two
%	designs differed only in which side the filled column fell on. They do not:
%	upstream's `sidebarleft` tints neither column and sets both in black on
%	white, so it has no filled column to mirror and now defines its own macros
%	in latexcv-sidebarleft.tex. Nothing else reads this file.
%
%	Upstream sets both as a pair of minipages side by side. A minipage cannot
%	break across a page, so that fixes the CV at one page and silently
%	overflows once it grows. `paracol` is used here instead: it flows two
%	independent columns across as many pages as they need, and can tint one of
%	them the whole way down.
%-------------------------------------------------------------------------------

\usepackage{paracol}

%-------------------------------------------------------------------------------
%	COLUMN PADDING
%-------------------------------------------------------------------------------
% Breathing space between the edges of the filled column and the text set on
% it. Upstream has none to inherit: its sidebar is a minipage of the full
% column width, and its contents sit hard against the tint on all four sides.
%
% The tint is grown outwards rather than the text inset, because the column is
% set centred -- a line inset by a parbox or by \leftskip would still centre on
% whatever it was inset into, and every centred heading in the column would
% shift with it. \backgroundcolor takes the extension, so the text column is
% left exactly where it is and only the colour reaches past it.
\newlength{\cvsidebarpad}
\setlength{\cvsidebarpad}{10pt}

% Gap between the filled column and the text beside it. The tint eats
% \cvsidebarpad of this on its way out, so the separation left between the two
% columns is what remains.
\setlength{\columnsep}{\dimexpr$if(columnsep)$$columnsep$$else$18pt$endif$+\cvsidebarpad\relax}

%-------------------------------------------------------------------------------
%	HEADINGS (main column)
%-------------------------------------------------------------------------------
\newcommand{\cvsection}[1]{%
  \vspace{4pt}
  \colorbox{sectcol}{\mystrut\makebox[\linewidth][l]{%
    \cvchevrons{bgcol}\hspace{4pt}%
    \textcolor{white}{\textbf{\uppercase{#1}}}\hspace{4pt}}}\\
  \vspace{2pt}}

\newcommand{\cvsubsection}[1]{%
  \vspace{4pt}
  {\larrow{sectcol}\hspace{2pt}\textcolor{bgcol}{\textbf{#1}}}\par
  \vspace{2pt}}

%-------------------------------------------------------------------------------
%	ENTRIES (main column)
%-------------------------------------------------------------------------------
% The main column is about two thirds of the page, so the date joins `with` at
% the right rather than taking a column of its own. The split is close to even
% because that right-hand cell carries two fields: at upstream's 0.55/0.42 an
% institution and a date range together spill onto a second line and leave the
% tail of the date stranded there on its own.
\newcommand{\cventryhead}[4]{%
  \begin{tabular*}{\linewidth}{@{}p{0.48\linewidth}@{\extracolsep{\fill}}x{0.48\linewidth}@{}}
    \textcolor{black}{\textbf{#2}} &
    \ifthenelse{\isempty{#3}}{}{\textcolor{complcol}{#3}}%
    \ifthenelse{\isempty{#4}}{}{\ifthenelse{\isempty{#3}}{}{, }\textit{\textcolor{complcol}{#4}}}%
    \ifthenelse{\isempty{#1}}{}{\ifthenelse{\isempty{#3}}{\ifthenelse{\isempty{#4}}{}{, }}{, }\textcolor{bgcol}{#1}}
  \end{tabular*}}

\newcommand{\vitaeentry}[5]{%
  \vspace{4pt}
  \cventryhead{#1}{#2}{#3}{#4}
  \vspace{-10pt}
  \textcolor{softcol}{\hrule}
  \vspace{6pt}
  #5}

\newcommand{\vitaebrief}[4]{%
  \vspace{4pt}
  \cventryhead{#1}{#2}{#3}{#4}}

% A skill over two lines, the category and then its members beneath it, rather
% than the two side by side as the other designs set them.
%
% This is upstream's own shape for a rated field: its `Fields` metasection sets
% the name on one line and the run of stars on the next. A rating is wide --
% ten stars are most of a column -- so beside a category it either crowds it or
% wraps, and the fold lands in the middle of the rating.
%
% The whole listing follows it, not only the rated entries, so that a listing
% mixing the two keeps one shape. \cvskillrow and the measured column it uses
% are still in the shared preamble for the five designs that set the two side
% by side; the measuring pass the listing runs is simply unused here.
%
% The two lines are one paragraph broken by \\*, not two paragraphs: \parskip
% is 6pt here, so a \par between them would set the members further from the
% category they belong to than the next entry is. Starred, so the pair cannot
% be split across a page. An empty second field is dropped rather than set as a
% blank line, so an entry with nothing but a category is one line.
\newcommand{\vitaeskill}[2]{%
  \vspace{2pt}
  \noindent\textbf{#1}%
  \ifthenelse{\isempty{#2}}{}{\\*#2}%
  \par\vspace{2pt}}

\newenvironment{vitaewhy}{%
  \begin{itemize}[label=\larrow{softcol}, topsep=0pt, partopsep=0pt, itemsep=3pt,
                  parsep=0pt, leftmargin=1.4em, labelsep=0.4em]%
}{%
  \end{itemize}%
}

%-------------------------------------------------------------------------------
%	SIDEBAR
%-------------------------------------------------------------------------------
% A heading in the filled column: white, uppercase, over a rule. Set as an
% environment rather than a command so the rule closes the group of items
% under it.
\newenvironment{cvsidebarsection}[1]{%
  \vspace{8pt}
  \begin{center}
    % A rating set on the filled column takes upstream's pair of colours for
    % it: the complementary accent for the filled stars, the column's own white
    % for the rest. The default pair is for a rating set on the page, and its
    % filled half is `sectcol` -- which is the colour of this column, so a
    % rating would otherwise disappear into it. Local to the environment, so
    % the main column is unaffected.
    \renewcommand{\cvstarfullcolor}{complcol}%
    \renewcommand{\cvstaremptycolor}{white}%
    % A bar is recoloured for the same reason and to the same pair: drawn in the
    % defaults its filled half is `sectcol` on `sectcol`, so the one part of it
    % that carries the rating is the one part that cannot be seen.
    \renewcommand{\cvbarfullcolor}{complcol}%
    \renewcommand{\cvbaremptycolor}{white}%
    % And on its own line, under whatever it rates, as upstream sets its
    % `Fields`. The column is a third of the page: ten stars are most of its
    % width, so a rating beside a label wraps and the fold lands mid-rating. A
    % bar takes the whole column once it has that line to itself.
    \renewcommand{\cvstarsbreak}{\\}%
    \renewcommand{\cvbarbreak}{\\}%
    \renewcommand{\cvbarwidth}{\linewidth}%
    \textcolor{white}{\large\uppercase{#1}}\\[2pt]
    \parbox{0.8\linewidth}{\textcolor{white}{\hrule}}\\[4pt]
    \normalsize
}{%
  \end{center}
  \vspace{-6pt}}

\newcommand{\cvsidebaritem}[1]{\textcolor{white}{#1}\\[2pt]}

% Contacts in the filled column: one per line, white, and centred, as upstream
% sets them -- its \icontext lines sit inside the `center` of a metasection,
% so each is centred on the column rather than aligned down a left edge.
% \begingroup keeps the redefinitions local: \cvcontact is the shared one
% again as soon as this returns, so a document emitting contacts twice does
% not get the sidebar's white, one-per-line form the second time.
\newcommand{\cvsidebarcontacts}{%
  \begingroup
  \cvcontactsreset
  \renewcommand{\cvcontactsep}{\\[3pt]}%
  \renewcommand{\cvcontact}[2]{%
    \ifcvfirstcontact\cvfirstcontactfalse\else\cvcontactsep\fi
    \textcolor{white}{\faIcon{##1}~##2}}%
  {\small\cvcontacts\par}%
  \endgroup}

%-------------------------------------------------------------------------------
%	TITLE (main column)
%-------------------------------------------------------------------------------
\newcommand{\cvtitleblock}{%
  \begin{flushleft}
    \HUGE\textsc{$name$$if(surname)$ $surname$$endif$}%
    $if(docname)$\cvtitlerule\textsc{$docname$}$endif$
    $if(position)$\\[4pt]\small\textsc{$position$}$endif$
  \end{flushleft}
  \vspace{2pt}
  \textcolor{softcol}{\hrule}
  \vspace{4pt}}

%-------------------------------------------------------------------------------
%	SIDEBAR CONTENT
%-------------------------------------------------------------------------------
% Filled from the document's `sidebar:` front matter, which is a list of
% sections each with a `title` and a list of `items`. The contact details are
% emitted first, since in these styles they belong in the column rather than
% under the name.
\newcommand{\cvsidebarbody}{%
  % Starred, because an ordinary \vspace at the top of a column is discarded --
  % which is exactly why the column's first heading sits hard against the tint
  % despite the one \cvsidebarsection opens with.
  \vspace*{\cvsidebarpad}%
$if(profilepic)$
  \begin{center}
    \includegraphics[width=$if(profilepic-width)$$profilepic-width$$else$0.7\linewidth$endif$]{$profilepic$}
  \end{center}
  \vspace{4pt}
$endif$
  \begin{cvsidebarsection}{$if(contacts-title)$$contacts-title$$else$Contact$endif$}
    \cvsidebarcontacts
  \end{cvsidebarsection}
$for(sidebar)$
  \begin{cvsidebarsection}{$sidebar.title$}
$for(sidebar.items)$
    \cvsidebaritem{$sidebar.items$}
$endfor$
  \end{cvsidebarsection}
$endfor$
$if(qrcode)$
  \begin{center}
    \includegraphics[width=$if(qrcode-width)$$qrcode-width$$else$0.55\linewidth$endif$]{$qrcode$}
  \end{center}
$endif$
}
